\section{Progressive exploration}\label{sec:ladders}

In this section we recall the framework of \emph{progressive exploration}, proposed by Fabia\'nski et al.~\cite{FabianskiPST19}, for  designing parameterized algorithms for \drds{r} and \dris{r}, among other problems of similar kind. Our main focus here is to explain that the progressive exploration algorithms for \drds{r} and \dris{r} can be executed using a bounded number of calls to the following two basic oracles, each working on a graph $G$ and stipulated by a distance parameter $r\in \N$:
\begin{itemize}
	\item $\far{S,r}$: Given a set $S$, return a vertex $v$ of $G$ such that $\dist_G(v,S)>r$; or $\bot$ if no such vertex exists.
	\item $\near{S,r}$: Given a set $S$, return a vertex $v$ of $G$ such that $\dist_G(v,s)\leq r$ for all $s\in S$; or $\bot$ if no such vertex exists.
\end{itemize}
This will allow us to concentrate in subsequent sections on implementing those oracles efficiently in the setting of dynamic data structures.

\subsection{Domination}
For \drds{r}, Fabia\'nski et al. proposed a \emph{semi-ladder algorithm} that proceeds as follows (see \cref{alg:1} for a pseudo-code). The algorithm is employed on a graph $G$ and searches for a distance-$r$ dominating set of size $k$, for given parameters $k,r\in \N$. It assumes access to the following two oracles, which will eventually be implemented using the more basic oracles $\far{S,r}$ and $\near{S,r}$:
\begin{itemize}
	\item $\notDominated{D}$:
	Given a set $D$ of $k$ vertices, return a vertex $v$ that is not distance-$r$ dominated by~$D$, that is, such that $\dist_G(v,D)>r$. In case no such vertex $v$ exists, return $\bot$.
	\item $\candidate{W}$: Given a vertex subset $W$, return a set $D$ of $k$ vertices that distance-$r$ dominates~$W$, that is, $\dist_G(w,D)\leq r$ for all $w\in W$. In case no such set $D$ exists, return $\bot$.
\end{itemize}
The algorithm iteratively constructs two sequences:
\begin{itemize}
	\item a sequence of \emph{candidates} $D_1,D_2,\ldots$, each being a set of at most $k$ vertices; and 
	\item a sequence of \emph{witnesses} $w_1,w_2,\ldots$, each being a single vertex of $G$.
\end{itemize}
After the $(i-1)$st iteration, candidates $D_1,\ldots,D_{i-1}$ and witnesses $w_1,\ldots,w_{i-1}$ are already constructed. Then the $i$th iteration proceeds as follows:
\begin{itemize}
	\item First, we call $\candidate{W}$ where $W$ comprises of all the  witnesses gathered so far, that is, $W\coloneqq \{w_1,\ldots,w_{i-1}\}$. This either yields the next candidate $D_i$ that at least distance-$r$ dominates $W$, or a conclusion that no such $D_i$ exists. In the latter case, we may conclude that $G$ has no distance-$r$ dominating set of size at most $k$.
	\item Second, we call $\notDominated{D_i}$ to verify whether $D_i$ already distance-$r$ dominates the whole graph. If so, then we have found a solution, and otherwise the call provides the next witness $w_i$, with which we may proceed to the next round of the iteration.  
\end{itemize}
It is clear that when this algorithm returns an answer, then this answer is always correct. The main insight of Fabia\'nski et al. is that since every candidate $D_i$ distance-$r$ dominates all the witnesses $w_1,\ldots,w_{i-1}$ but not the witness $w_i$, the constructed candidates and witnesses form a pattern called a \emph{semi-ladder}. As such semi-ladders cannot be too long in classes of bounded expansion, this provides a bound on the number of iterations that the algorithm executes before breaking the loop and reporting an outcome.


\begin{algorithm}
	\SetKwInOut{Input}{Input}
	\SetKwInOut{Output}{Output}
	
	\nonl
	\underline{method DomSet} $(G,r,k)$\\
	\Input{A graph $G$ and parameters $r,k\in \N$}
	\Output{A distance-$r$ dominating set of size $k$, or $\bot$ if no such set exists.}
	\medskip
	$W \coloneqq \emptyset$\;
	\Repeat
	{
		$D\coloneqq \candidate{W}$\;
		\If{$S = \bot$}
		{
			\KwRet{$\bot$} \;
		}
		$w \coloneqq \notDominated{D}$ \;
		\If{$w = \bot$}
		{
			\KwRet{$S$} \;
		}
		$W \coloneqq W \cup \{ w \}$ \;
	}
	\caption{Semi-ladder algorithm for distance-$r$ dominating set of size $k$ in a graph $G$}\label{alg:1}
\end{algorithm}

To be more precise, a \emph{distance-$r$ semi-ladder} of order $\ell$ in a graph $G$ is a pair of sequences of vertices $a_1,\ldots,a_\ell$ and $b_1,\ldots,b_\ell$ such that
\begin{itemize}
	\item $\dist_G(a_i,b_i)>r$ for each $i\in [\ell]$; and
	\item $\dist_G(a_i,b_j)\leq r$ for all $1\leq j<i\leq \ell$.
\end{itemize}
The \emph{distance-$r$ semi-ladder index} of a graph $G$, denoted $\sli_r(G)$, is the largest order of a distance-$r$ semi-ladder that can be found in $G$. For a graph class $\Cc$, we define its distance-$r$ semi-ladder index as $\sli_r(\Cc)\coloneqq \sup_{G\in \Cc} \sli_r(G)$. Note that this value may be infinite if graphs from $\Cc$ contain distance-$r$ semi-ladders of arbitrarily large order. But as proved by Fabia\'nski et al., this does not happen in bounded expansion~classes.

\begin{theorem}[{\cite[Lemma~29]{FabianskiPST19}}]\label{thm:sli-be}
	For every class $\Cc$ of bounded expansion and $r\in \N$, $\sli_r(\Cc)$ is finite.
\end{theorem}

We remark that \cref{thm:sli-be} holds even in a larger generality when $\Cc$ is nowhere dense, but we will not use this here. We also note that for $r=1$, excluding a biclique suffices to bound the semi-ladder index.

\begin{theorem}[{\cite[Lemma~33]{FabianskiPST19}}]\label{thm:sli-biclique}
	For every graph $G$ that does not contain the biclique $K_{t,t}$ as a subgraph, for some $t\in \N$, we have $\sli_1(G)<3t$.
\end{theorem}

The following statement summarizes how a bound on the semi-ladder index influences the number of iterations executed by the semi-ladder algorithm.

\begin{theorem}[consequence of {\cite[Lemma~5 and Lemma~24]{FabianskiPST19}}]\label{thm:termination}
	Let $\Cc$ be a graph class and $r\in \N$ be such that $\sli_r(\Cc)$ is finite. Then the semi-ladder algorithm (\cref{alg:1}) run on any graph from $\Cc$ terminates after executing less than $k^{\sli_r(\Cc)}$ iterations.
\end{theorem}

Note here that if the algorithm performs at most $\ell\coloneqq k^{\sli_r(\Cc)}$ iterations in total, then all the calls to the $\candidate{W}$ oracle are applied only to sets $W$ of size at most $\ell$.

Our goal now is to reduce implementation of the $\notDominated{D}$ and $\candidate{W}$ oracles to the more basic oracles $\far{S,r}$ and $\near{S,r}$. For $\notDominated{D}$ this is trivial: queries $\notDominated{D}$ and $\far{D,r}$ are just equivalent. For $\candidate{W}$, this is a bit more complex.

\begin{lemma}\label{lem:candidate-to-near}
	Let $G$ be a graph, $W$ be a subset of vertices of $G$, and $k,r\in \N$ be parameters. Then the query $\candidate{W}$ can be answered in time $\Oh(k^{|W|+1})$ by performing at most $k^{|W|+1}$ queries of the form $\near{S,r}$, where $S$ is a subset of $W$.
\end{lemma}
\begin{proof}
	We enumerate all
	the partitions of $W$ into at most $k$ subsets; there are at most $k^{|W|}$ such partitions. For every partition $\cal P$ and every part $S\in \cal P$, we issue the query $\near{S,r}$ to verify whether there exists a vertex $v_S$ such that $\dist_G(v_S,s)\leq r$ for all $s\in S$. If for some partition $\cal P$ we can find such a vertex $v_S$ for every part $S\in \cal P$, then the set $D\coloneqq \{v_S\colon S\in \cal P\}$ is a valid answer to the query $\candidate{W,r}$. And if this check fails for every considered partition $\cal P$, then there is no set $D$ of size at most $k$ that distance-$r$ dominates $W$ and we can safely return $\bot$.
	Thus, the total
	number of queries $\near{S,r}$ issued is bounded by
	$k^{|W|+1}$.
\end{proof}

We conclude this section by a statement summarizing the discussion and implying the following: designing an efficient data structure for the $\near{S}$ and $\far{S}$ queries suffices to implement the semi-ladder algorithm in the dynamic setting, thereby providing a dynamic data structure for the \drds{r} problem.

\begin{lemma}
	Let $\Cc$ be a class of graphs, $k,r\in \N$ be such that $\sli_r(\Cc)$ is finite, and $\ell\coloneqq k^{\sli_r(\Cc)}$. Let $G$ be a dynamic graph that belongs to $\Cc$ at all times. Suppose that there are data structures $\nearVertexDS{G}{r,\ell}$, supporting queries $\near{S,r}$ in $G$ with $|S|\leq \ell$, and $\farVertexDS{G}{r,k}$, supporting queries $\far{S,r}$ in $G$ with $|S|\leq k$, so that $\nearVertexDS{G}{r,\ell}$ and $\farVertexDS{G}{r,k}$ have amortized update/query time $T$, initialization time $I$, and space complexity $M$. Then there is a data structure $\dominatingSetDS{G}{r,k}$ for $G$ that supports the~query
	\begin{itemize}
		\item $\ds{}$: Return a distance-$r$ dominating set of size at most $k$ in $G$, or $\bot$ if no such set exists.
	\end{itemize}
	The amortized update time of $\dominatingSetDS{G}{r,k}$ is $2^{k^{\Oh(\sli_r(\Cc))}}\cdot T$, the query time is $\Oh(k)$, the initialization time is $\Oh(I)$, and the space complexity is $\Oh(M)$.
	
	Moreover, the data structures $\nearVertexDS{G}{r,\ell}$ and $\farVertexDS{G}{r,k}$ may be randomized with error probability at most $1-\eps$, for any fixed parameter $\eps>0$ and against an oblivious adversary, with amortized update/query time, initialization time, and space complexity becoming $T\cdot \tfrac{1}{\eps}$, $I\cdot \log \tfrac{1}{\eps}$, and $M\cdot \log \tfrac{1}{\eps}$, respectively. In this case,  $\dominatingSetDS{G}{r,k}$ is also randomized with error probability at most $\eps$ against an oblivious adversary. The amortized update time becomes $2^{k^{\Oh(\sli_r(\Cc))}}\cdot T\cdot \tfrac{1}{\eps}$, the initialization time becomes $k^{\Oh(\sli_r(\Cc))}\cdot I\cdot \log \tfrac{1}{\eps}$, and the space complexity becomes  $k^{\Oh(\sli_r(\Cc))}\cdot M\cdot \log \tfrac{1}{\eps}$.
\end{lemma}
\begin{proof}
	The data structure $\dominatingSetDS{G}{r,k}$ simply maintains the assumed data structures $\nearVertexDS{G}{r,\ell}$ and $\farVertexDS{G}{r,k}$. Upon every update to $G$, we relay the update to $\nearVertexDS{G}{r,\ell}$ and $\farVertexDS{G}{r,k}$ and we run the semi-ladder algorithm on $G$ to find a distance-$r$ dominating set of size at most $k$ (or detect a lack thereof). By \cref{thm:termination}, this algorithm can be implemented using at most $k^{\sli_r(\Cc)}$ queries $\notDominated{D}$ oracle with $|D|\leq k$ and $k^{\sli_r(\Cc)}$ queries $\candidate{W}$ with $|W|\leq \ell$. Every query $\notDominated{D}$ can be answered by a single call to the query $\far{D,r}$ offered by $\farVertexDS{G}{r,k}$. By \cref{lem:candidate-to-near}, every query $\candidate{W}$ can be answered by performing $k^{|W|+1}\leq 2^{k^{\Oh(\sli_r(\Cc))}}$ calls to queries of the form $\near{S,r}$ with $S\subseteq W$, offered by $\nearVertexDS{G}{r,\ell}$. Once a distance-$d$ dominating set of size $k$ (or lack thereof) is computed, it can be provided in time $\Oh(k)$ upon any $\ds{}$ update.
	The claimed bounds on the amortized update time of $\dominatingSetDS{G}{r,k}$ follow immediately from the assumed guarantees about $\nearVertexDS{G}{r,\ell}$ and $\farVertexDS{G}{r,k}$; and similarly for the initialization time and the space complexity.
	
	In case $\nearVertexDS{G}{r,\ell}$ and $\farVertexDS{G}{r,k}$ are randomized, $\dominatingSetDS{G}{r,k}$ maintains their instances with a rescaled error parameter $\eps'\coloneqq \tfrac{\eps}{q}$, where
	\[q\coloneqq k^{\sli_r(\Cc)}+k^{k^{\sli_r(\Cc)}}\leq 2^{k^{\Oh(\sli_r(\Cc))}}\]
	is an upper bound on the total number of queries to $\nearVertexDS{G}{r,\ell}$ or $\farVertexDS{G}{r,k}$ that may occur in a single run of the semi-ladder algorithm. Therefore, by the union bound, we conclude that with probability at least $1-\eps$ none of these queries returns an incorrect answer and the semi-ladder algorithm runs correctly. It is straightforward to verify that rescaling the error probability in $\nearVertexDS{G}{r,\ell}$ and $\farVertexDS{G}{r,k}$ influences the complexity guarantees for $\dominatingSetDS{G}{r,k}$ as claimed.
\end{proof}

%The algorithm to compute $r$-domination set of size $k$ is then implemented using these procedures
%and the technique called by the authors \emph{progressive exploration}.
%The progressive exploration technique constructs a sequence of candidate solutions. For each
%candidate solution that it constructs, it then tries to invalidate it by finding a witness vertex
%that is not dominated by it. Then the next candidate solution is constructed to dominate both the
%current witness and previously explored witnesses. This process is shown in
%Algorithm~\ref{alg:1}.


%To obtain a bound on the running time of this simple algorithm the
%authors introduce a number of tools from logic and combinatorial graph theory.
%In particular, they define a bipartite graph induced by what they call
%a distance formula. This fomula is a logical sentence stating
%some specific property of a graph $G$ that is an input instance to the
%$r$-dominating set problrem.
%For the $r$-dominating set problem, the
%distance formula is
%$\delta^k_r(\overline{x},y)=\bigvee_{i=1}^k \delta_r(x_i,y)$,
%stating that there is some $x_i$ in a touple
%$\overline{x}=(x_1,\ldots,x_k)$ such that the distance from $x_i$ to $y$
%is at most $r$. This is a First Order Logic formula, whose length
%depends linearly on two parameters: $k$-the size of $\overline{x}$ and
%$r$ - the distance. The bipartite graph $B(\delta^k_r,G)$ induced
%by formula
%$\delta^k_r$ on a graph $G$ is then defined as follows. The left side
%$L$ consists of all tuples $(v_1,\ldots,v_k) \subseteq V^k$ of
%vertices of $G$. The left side $L$ represents
%the universum of all candidate solutions for the $r$-dominating
%set that we seek. The right side $R$ consists of
%all vertices $v \in V(G)$. The right side is viewed as a set of
%witnesses. There is an edge between a candidate
%$(v_1,\ldots,v_k) \in L$ and a witness $v \in R$ if and only if
%$\delta^k_r((v_1,\ldots,v_k),v)$ holds, i.e., $v$ is at
%distance $r$ from some vertex in $(v_1,\ldots,v_k)$. Clearly,
%a feasible solution is a neighbour in $B(\delta^k_r,G)$ to all
%vertices $v \in R$. Unfortunately we cannot afford to construct
%the entire graph $B(\delta^k_r,G)$.
%Algorithm~\ref{alg:1} searches $B(\delta^k_r,G)$ by successively
%exploring new candidates and witnesses, one candidate and one
%witness in each round of the while loop.
%It is easy to observe (see~\cite{FabianskiPST19}), that the
%candidates and witnesses explored by the algorithm induce a
%semi-ladder in $B(\delta^k_r,G)$. Thus, the authors define the
%semi-ladder index $\sli(B)$ of a bipartite graph $B$ as the size of the
%largest induced semi-ladder contained in $B$. Thus, the number of
%rounds executed by Algorithm~\ref{alg:1} is bounded by
%$\sli(B(\delta^k_r,G))$. This in turn, as we shall see later,
%is nicely bounded from above
%for some specific graph classes, in particular bounded expansion graphs.

%By Theorem 11 and Corollary 10 in \cite{FabianskiPST19}, we get that on nowhere dense graphs
%the progressive exploration algorithm (Algorithm~\ref{alg:1})
%performs a bounded number rounds, since the number of rounds corresponds to the size
%of a semi-ladder in a bipartite graph induced by a distance formula
%(see \cite{FabianskiPST19} for the details).

\subsection{Independence}

\michalin{Work in progress}


\begin{algorithm}
	\SetKwInOut{Input}{Input}
	\SetKwInOut{Output}{Output}
	
	\nonl
	\underline{method IndSet} $(G,r,k,p)$\\
	\Input{A graph $G$ and parameters $r,k,p\in \N$}
	\Output{A distance-$r$ independent set of size $k$, or $\bot$ if no such set exists.}
	\medskip
	${\cal L}\coloneqq \emptyset$\;
	$W \coloneqq \emptyset$\;
	\Repeat
	{
		$I\coloneqq \candidateInd{W}$\;
		\If{$I = \bot$}
		{
			\KwRet{$\bot$} \;
		}
		${\cal L}\coloneqq {\cal L}\cup \{I\}$ \;
		$P \coloneqq \witnessInd{\cal L,p}$ \;
		\If{$P = \bot$}
		{
			\KwRet{$I$} \;
		}
		$W \coloneqq W \cup P$ \;
	}
	\caption{Ladder algorithm for distance-$r$ independent set of size $k$ in a graph $G$}\label{alg:2}
\end{algorithm}


\subsection{Old text}
In this work, we use the above framework to implement efficient
dynamic algorithms for the $r$-dominating set problem.
We explore the fact that we have the bound on the number of rounds
of Algorithm~\ref{alg:1} at hand. What remains to show is how
to efficiently answer $\candidate{W}$ and $\notDominated{S}$
queries in a dynamically changing graph. For that, we introduce
two dynamic data structures that are the building blocks for the above.

\begin{definition}
The $\nearVertexDS{G}{k,d}$ data structure maintains a dynamic graph $G$ by allowing the following operations:
\begin{itemize}
 \item $\init{G}$ - initializes with a graph $G$ on $n$ vertices
 \item $\add{uv}$ - inserts an edge $uv$ to $G$
 \item $\remove{uv}$ - removes an edge $uv$ from $G$
 \item $\near{S}$ - returns a vertex at distance at most $r$ from all vertices in $S$ if such vertex exists, otherwise it returns $\bot$.
\end{itemize}
\end{definition}

\begin{definition}
The $\farVertexDS{G}{k,d}$ data structure maintains a dynamic graph $G$ by allowing the following operations:
\begin{itemize}
 \item $\init{G}$ - initializes with a graph $G$ on $n$ vertices
 \item $\add{uv}$ - inserts an edge $uv$ to $G$
 \item $\remove{uv}$ - removes an edge $uv$ from $G$
 \item $\far{S}$ - returns a vertex at distance at least $r+1$ from
  $S$ if such vertex exists, otherwise it returns $\bot$.
\end{itemize}
\end{definition}


\begin{theorem}
Using $\nearVertexDS{G}{k,d}$ and $\farVertexDS{G}{k,d}$ data
structures,
one can implement Algorithm~\ref{alg:1} in time
$\sli(B(\delta^k_r,G)) \cdot ( k^{\sli(B(\delta^k_r,G))+1}
\mathsf{Q_{nv}} +
2k \cdot \mathsf{U_{fv}} + \mathsf{Q_{fv}})$, where
$\mathsf{Q_{nv}}$ is the query time
of $\nearVertexDS{G}{k,d}$ and $\mathsf{U_{fv}}$ and
$\mathsf{Q_{fv}}$ is the update
and the query time of
$\farVertexDS{G}{k,d}$ respectively.
\end{theorem}
\begin{proof}We show how to implement procedures $\candidate{W}$ and $\notDominated{S}$.
\paragraph*{Candidate.}We set up and maintain a
$\nearVertexDS{G}{k,d}$ data structure on the
input dynamic graph $G$. We implement query $\candidate{W}$
as follows.
We enumerate all
the partitions of vertices in $W=W_1 \cup \ldots \cup W_k$
into $k$ subsets. There is at most $k^{|W|}$ such partitions.
For each
such partition and each $i \in [k]$ we issue $v_i=\near{W_i}$
query, that returns a vertex at distance at most $r$ from
all vertices in $W_i$. If for some partition we can find near
vertices for all $i \in [k]$, then the set
$S=\{v_1,\ldots,v_k \}$ is a feasible candidate and we return
it. Otherwise, we return $\bot$. Thus, the total
number of queries issued to $\nearVertexDS{G}{k,d}$ is
$k^{|W|+1}$.

\paragraph*{Witness.} We set up and maintain a
$\farVertexDS{G}{k,d}$ data
structure on the input dynamic graph $G$ and some auxiliary
vertex $q \notin V(G)$. Each time a query
$\notDominated{S}$ is used by Algorithm~\ref{alg:1},
we issue $\far{S}$ query, which returns a vertex $v$ at
distance at least $r+1$ from $S$. Thus, $v$ is
a valid witness that $S$ is not a feasible solution, or, in
other words $\delta^k_r(S,v)$ does not hold.
We return a vertex $v$ if it exists,
or $\bot$ otherwise. This takes time proportional
to a query operation on $\farVertexDS{G}{k,d}$ data structure.

\end{proof}

The first step to use the above theorem is to bound the ladder index
$\sli(B(\delta_r^k,G))$ for specific classes of graphs.
After~\cite{FabianskiPST19} we get the following result.

\begin{theorem}[Theorem 11 and Lemma 5 in ~\cite{FabianskiPST19}]
For any $r,k \in \N$ and a nowhere dense class $\mathcal{C}$ of graphs
it holds that
$\sli(B(\delta_r^k,G)) \in k^{\Oh_{\mathcal{C},r}(1)}$ for any
$G \in \mathcal{C}$.
\end{theorem}

In Section~\ref{sec:distone} we develop the
following data structures for degenerate graphs.

\begin{theorem}\label{thm:distone-near}
Assume $n\ge 2$ and $d\ge 1$.
There exists
$\nearVertexDS{G}{k,d}$ data structure which
on a dynamic $d$-degenerate graph $G$ achieves the
following running times:
\begin{itemize}
\item $\init{G}$:  $\Oh(nd\cdot 2^{4d}(d^3+\log^3 n))$ time;
\item $\add{uv}$ and $\remove{uv}$:
$\Oh(2^{4d}(d^3+\log^3 n))$ amortized time;
\item $\near{S}$: $\Oh(|S| + d|S|^2\log n)$ time.
\end{itemize}
The structure uses $\Oh(nd\cdot 2^{4d})$ words of space.
\end{theorem}

\begin{theorem}\label{thm:distone-far}
Assume $n\ge 2$, $d\ge 1$, and let $\eps \in (0,1/2)$.
There exists a randomized data structure $\farVertexDS{G}{k,d}$ that,
against an oblivious adversary, maintains a dynamic $d$-degenerate
graph $G$ achieving the following running times:
\begin{itemize}
\item $\init{G}$:
$\Oh\!\left(nd\cdot 2^{4d}(d^3+\log^3 n)\log n \log \frac{1}{\eps}\right)$;
\item $\add{uv}$ and $\remove{uv}$: amortized
$\Oh\!\left(2^{4d}(d^3+\log^3 n)\log n \log \frac{1}{\eps}\right)$ time;
\item $\far{S}$: admits error probability bounded by $\epsilon$, takes time
\[
\Oh\!\left((|S|+2^{|S|} d {|S|}^2\log n)\log n \log \frac{1}{\eps}\right).
\]
\end{itemize}
The structure uses $\Oh\!\left(nd\cdot 2^{4d}\log n \log \frac{1}{\eps}\right)$ words of space.
\end{theorem}

By \cite{grohe_et_al:LIPIcs.FSTTCS.2013.21} bounded expansion graphs
are both nowhere dense and have bounded degeneracy. As a corollary, we
obtain the following result.


\begin{corollary}
Let $\mathcal{C}$ be a bounded expansion class of graphs.
There exists a fully dynamic $\dominatingSetDS{G}{r,k}$ data structure,
that maintains a graph $G \in \mathcal{C}$ by supporting the following
operations:
\begin{itemize}
\item $\init{G}$: initializas the data structure in time
\item $\add{uv}$ and $\remove{uv}$:
\item $\ds{k}$:
\end{itemize}

TODO: details.
\end{corollary}

\azinline{i need section 5 to write the conclusions}

\begin{theorem}\label{thm:NearVertex}
There is a $\nearVertexDS{G}{k,d}$ data structure that maintains a dynamic graph $G$ of bounded
expansion, providing the following operations:
\begin{itemize}
 \item $\init{G}$ - initializes with a graph $G$ on $n$ vertices
 \item $\add{uv}$ - inserts an edge $uv$ to $G$
 \item $\remove{uv}$ - removes an edge $uv$ from $G$
 \item $\near{S}$ - returns a vertex at distance at most $r$ from all vertices in $S$ if such vertex exists, otherwise it returns $\bot$.
\end{itemize}
\end{theorem}
\begin{proof}Postponed to Section 5. \end{proof}

\begin{theorem}\label{thm:FarVertex}
There is a $\farVertexDS{G}{k,d}$ data structure that maintains a dynamic graph $G$ of bounded
expansion, providing the following operations:
\begin{itemize}
 \item $\init{G}$ - initializes with a graph $G$ on $n$ vertices
 \item $\add{uv}$ - inserts an edge $uv$ to $G$
 \item $\remove{uv}$- removes an edge $uv$ from $G$
 \item $\far{v}$ - returns a vertex at distance at least $r$ from $v$ if it exists,
 otherwise it returns $\bot$
\end{itemize}
\end{theorem}
\begin{proof}Postponed to Section 5. \end{proof}


